\documentclass[10pt]{article}

\usepackage[margin=1in]{geometry}
\usepackage{booktabs}
\usepackage{hyperref}

\title{mcp-stdio-guard: Protocol Hygiene Evaluation for MCP Stdio Servers}
\author{}
\date{}

\begin{document}
\maketitle

\begin{abstract}
This paper evaluates MCP stdio server failure modes around stdout protocol
hygiene, JSON-RPC framing, lifecycle conformance, process behavior, and
registry-facing result contracts. Experimental results in this paper must be
generated from the evaluation scripts in this artifact repository and tied to a
pinned mcp-stdio-guard release or commit.
\end{abstract}

\section{Introduction}

MCP stdio transports reserve stdout for JSON-RPC protocol messages. Diagnostic
output, startup banners, invalid frames, missing lifecycle responses, and
process failures can therefore surface as confusing client behavior. This
artifact separates the research evaluation from the production
\texttt{mcp-stdio-guard} repository.

\section{Methodology}

All experiments should invoke a pinned \texttt{mcp-stdio-guard} version against
synthetic and, when appropriate, real MCP stdio servers. Generated artifacts
should record the guard version, fixture command, exit status, and raw JSON
output.

\section{Results}

Results are intentionally omitted from the source tree until generated by an
actual evaluation run.

\section{Artifact Availability}

The production tool is available at
\url{https://github.com/1Utkarsh1/mcp-stdio-guard}. This repository contains
the research artifact and reproducibility notes.

\bibliographystyle{plain}
\bibliography{references}

\end{document}
